\chapter*{结语}\phantomsection\addcontentsline{toc}{section}{结语}

回到本文开头的问题：怎样让客户机拥有一台机器的使用接口，却不拥有对真实资源的最终控制权。三种实现路径给出的答案并不构成一条简单的替代链，它们是同一问题在不同约束下的三种实现；其中嵌套一并两章，因为“把多层语义压到单层硬件上”与“削减由此产生的跨层往返”是两个先后相接、判据不同的问题。沿四维框架收束，可以得出三点判断。

第一，四个维度必须协作，任何单一维度都无法独立构成一台可用的虚拟机。CPU、内存、中断、设备各自解决执行、地址、异步事件与数据通路，而这些状态必须相互一致：地址转换的结果决定了访存与DMA是否落到正确的内存，中断投递要依赖CPU的执行上下文与设备的后端状态，设备完成通知又回到中断。软件方案里四者靠同一个进程内的数据结构连接，硬件方案里它们靠同一套异常级与系统结构的分工连接，嵌套方案里它们靠同一组跨层映射与事件归属连接\cite{source3,source4,source7}。割裂地看任何一个维度，都看不到整台机器为何能够运行。

第二，硬件辅助减少的是高频路径上的软件介入，管理策略与设备语义仍然需要软件。Stage-2把地址翻译下放给硬件，但内存后备的分配与映射维护仍是KVM的工作；在相应配置下，GICv4能向驻留vPE直接注入vLPI；映射配置、调度和doorbell响应仍需软件参与；SMMU约束真实DMA，却不能替代设备模型去实现网卡或磁盘的功能\cite{source4,source9,source10}。硬件提供的是机制，把机制组合成服务并决定何时使用它，仍是软件的责任。

第三，嵌套虚拟化的本质是把资源边界与上层策略组合起来，而这一步分成两个问题。Turtles 展示了将多层语义组合到现有硬件支持上的实现：VMX 指令由 L0 代执行、VMCS02 由两份配置合成、多层地址翻译折叠进一张 EPT，代价则集中在跨层往返与 L1 处理时的逐级放大~\cite{source7}。HyperTurtle 接着回答第二问，其适用范围是被安全下放且一致性成本可控的热路径——它依赖 verifier、受限 helper 与共享状态的同步协议维持相应隔离边界，并以回退路径兜底，因此优化的是适用的热路径，不取消 L1~\cite{source8}。两章合起来才构成完整判断：先把嵌套做对，再谈削减它的成本；前者的收益不能解释后者的加速比。

最后需要说明证据的边界。HyperTurtle的评测在固定绑核、关闭SMT、\codeplain{idle=poll}与\codeplain{nopti}的收紧配置下取得，网络吞吐实验还调整了vCPU数量，这些配置限定了实验中的资源争用、空闲等待和隔离开销条件\cite{source8}。因此论文报告的相对收益说明的是该配置下热路径优化的效果，不能直接外推为共享多租户生产环境的期望收益；本文未进行性能复现，对更广泛部署的适用性仍需通过独立实验检验。

把三种实现并置（嵌套两章合为一层来看），真正可比较的是执行、状态与策略被放在了哪里：软件把三者都放在进程内，硬件把机制下放而保留管理，嵌套则在资源边界与上层策略之间重新分配处理逻辑。虚拟化技术的演进，最终是在减少跨层操作与维持控制权之间不断寻找新的平衡点。

下表把三种方案在同一四维框架下并置，作为全文的收束。

\begin{table}[htbp]
\centering
\small
\begin{tabularx}{\linewidth}{@{}lYYY@{}}
\toprule
维度 & QEMU软件 & Arm硬件辅助 & x86-64嵌套 \\
\midrule
CPU & 翻译执行与软件CPU状态，guest指令由宿主代码实现 & guest在EL1/EL0原生执行，EL2控制，KVM切换上下文 & 虚拟化VMX与合成VMCS，按退出归属决定就地处理或转发 \\
内存 & 软件MMU与TLB，解释guest页表并映射到宿主内存 & Stage-1与Stage-2两阶段翻译，KVM管理Stage-2映射与相关异常 & 合成EPT压缩多层映射，维护映射一致性并下放部分缺页处理 \\
中断 & 软件控制器维护pending，经架构异常投递到guest & GICv3虚拟接口安排注入，GICv4对驻留vPE直接注入vLPI & 跨层事件转发与注入，减少特定处理路径的往返 \\
设备 & MMIO设备模型与软件DMA，后端完成请求并通知guest & 模拟、virtio与直通三种形态并存，SMMU约束真实DMA & 设备与IOMMU映射组合，下放数据处理与上层策略 \\
\bottomrule
\end{tabularx}
\caption{四维框架下的三种实现对照}
\label{tab:conclusion-dims}
\end{table}
